% Work-MAT | 2. rocnik | hodina 14
\documentclass[11pt,a4paper]{article}
\input{../../sablony/pripravy/priprava.sty}
\nastavPripravu{2. ročník}{Exponenciální funkce}{14}{Sestavování předpisu exponenciální funkce}
\begin{document}
\poznamka{Navazuje na tvorbu grafu a jeho posuny. Dnešní otázka je obrácená: místo kreslení z předpisu budou žáci sestavovat předpis z několika údajů. Zatím nepoužívat logaritmy ani exponenciální rovnice s neznámým exponentem.}
\blok{Z bodu zpět k předpisu}
Základní exponenciální funkci popisuje předpis $f(x)=a^x$. Známe-li bod grafu s nenulovou první souřadnicí, můžeme hledat základ $a$. Při určování odmocniny vždy respektujeme podmínku $a>0$.
\begin{prikladReseniBox}[title={Příklad 1}]
Urči předpis $f(x)=a^x$, pokud graf prochází bodem $[2;9]$. Sestroj ještě dva snadno počitatelné body grafu.
\tcblower
\[
\begin{aligned}
f(2)&=9,\\
a^2&=9,\\
a&=3\quad(a>0).
\end{aligned}
\]
Proto $f(x)=3^x$. Na grafu leží například body $[0;1]$ a $[1;3]$.
\end{prikladReseniBox}
\begin{prikladReseniBox}[title={Příklad 2}]
Graf základní exponenciální funkce prochází bodem $[-1;4]$. Sestav její předpis a urči monotónnost.
\tcblower
\[
\begin{aligned}
a^{-1}&=4,\\
\frac{1}{a}&=4,\\
a&=\frac{1}{4}.
\end{aligned}
\]
Předpis je $f(x)=\left(\frac{1}{4}\right)^x$ a funkce klesá.
\end{prikladReseniBox}
\blok{Když funkce nezačíná na hodnotě jedna}
V aplikacích často známe počáteční hodnotu $b$ a za každý další krok násobíme stálým číslem $a$. Pak používáme model $g(x)=b\cdot a^x$, kde $b>0$ a $a>0$, $a\ne1$. Nejde o základní tvar $a^x$, ale o jeho svisle násobenou variantu.
\begin{prikladReseniBox}[title={Příklad 3}]
Sestav předpis $g(x)=b\cdot a^x$, jestliže $g(0)=5$ a $g(1)=15$.
\tcblower
\[
\begin{aligned}
g(0)&=b\cdot a^0=b=5,\\
g(1)&=5a=15,\\
a&=3.
\end{aligned}
\]
Hledaná funkce je $g(x)=5\cdot3^x$. Kontrola: $g(0)=5$, $g(1)=15$.
\end{prikladReseniBox}
\Needspace{10\baselineskip}
\begin{prikladReseniBox}[title={Příklad 4}]
Funkce $h(x)=b\cdot a^x$ splňuje $h(0)=2$ a $h(2)=18$. Urči předpis a hodnotu $h(3)$.
\tcblower
\[
\begin{aligned}
h(0)&=b=2,\\
h(2)&=2a^2=18,\\
a^2&=9,\\
a&=3\quad(a>0).
\end{aligned}
\]
Tedy $h(x)=2\cdot3^x$ a $h(3)=54$.
\end{prikladReseniBox}
\poznamka{Dobrý rychlý test porozumění: v předpisu $b\cdot a^x$ určujeme $b$ z hodnoty pro $x=0$; teprve potom hledáme $a$. Srovnej s lineární funkcí, kde se po jednotkovém kroku hodnoty sčítají.}
\begin{zadaniReseniBox}[title={Samostatná práce 1}]
\begin{enumerate}[label=\alph*)]
\item Urči $f(x)=a^x$, pokud $f(2)=\frac{1}{25}$.
\item Urči $g(x)=b\cdot a^x$, pokud $g(0)=3$ a $g(1)=12$.
\item Urči $h(x)=b\cdot a^x$, pokud $h(0)=8$ a $h(2)=2$.
\end{enumerate}
\tcblower
\begin{enumerate}[label=\alph*)]
\item $a^2=\frac{1}{25}$, tedy $a=\frac{1}{5}$ a $f(x)=\left(\frac{1}{5}\right)^x$.
\item $b=3$ a $3a=12$, tedy $g(x)=3\cdot4^x$.
\item $b=8$, $8a^2=2$, $a=\frac{1}{2}$, tedy $h(x)=8\cdot\left(\frac{1}{2}\right)^x$.
\end{enumerate}
\end{zadaniReseniBox}
\begin{zadaniReseniBox}[title={Rozšíření 1}]
Předpis $k(x)=a^{x-p}+2$ má asymptotu $y=2$ a prochází body $[1;3]$ a $[2;5]$. Urči $a>0$ a $p$, víš-li, že $a\ne1$.
\tcblower
Z rovností $a^{1-p}=1$ a $a^{2-p}=3$ plyne $1-p=0$, tedy $p=1$; druhá rovnost pak dává $a=3$. Proto $k(x)=3^{x-1}+2$.
\end{zadaniReseniBox}
\begin{zadaniReseniBox}[title={Domácí úkol}]
\begin{enumerate}[label=\alph*)]
\item Sestav $f(x)=a^x$ z bodu $[-1;2]$.
\item Sestav $g(x)=b\cdot a^x$ z hodnot $g(0)=4$, $g(1)=6$.
\item Vysvětli, jaký význam mají parametry $a$ a $b$ v předpisu $b\cdot a^x$.
\end{enumerate}
\tcblower
\begin{enumerate}[label=\alph*)]
\item $a^{-1}=2$, tedy $f(x)=\left(\frac{1}{2}\right)^x$.
\item $b=4$, $4a=6$, $a=\frac{3}{2}$, tedy $g(x)=4\cdot\left(\frac{3}{2}\right)^x$.
\item $b$ je počáteční hodnota a $a$ je násobný koeficient při zvýšení $x$ o jednu.
\end{enumerate}
\end{zadaniReseniBox}
\end{document}
